\documentclass[11pt,a4paper]{amsart}
\usepackage[T1]{fontenc}
\usepackage{lmodern}
\usepackage{amssymb}
\usepackage{microtype}
\usepackage[hyphens]{url}
\usepackage[colorlinks=true,linkcolor=black,citecolor=black,urlcolor=black]{hyperref}
\usepackage[margin=2.5cm]{geometry}
\usepackage{tikz}
\usetikzlibrary{shapes.geometric}

\hypersetup{
  pdftitle={Real quadratic planes that need four colours},
  pdfauthor={Sergi Gal\'an},
  pdfsubject={The Hadwiger-Nelson problem over number fields},
  pdfkeywords={Hadwiger-Nelson problem, chromatic number of the plane, unit-distance graph, real quadratic field, SAT, DRAT}}

\newtheorem{theorem}{Theorem}
\newtheorem{corollary}[theorem]{Corollary}
\newtheorem{proposition}[theorem]{Proposition}
\newtheorem{lemma}[theorem]{Lemma}
\theoremstyle{remark}
\newtheorem*{remark}{Remark}

\newcommand{\F}{\mathbb{F}}
\newcommand{\Q}{\mathbb{Q}}
\newcommand{\R}{\mathbb{R}}
\newcommand{\Z}{\mathbb{Z}}
\newcommand{\C}{\mathbb{C}}
\newcommand{\p}{\mathfrak{p}}
\newcommand{\PP}{\mathfrak{P}}

\title{Real quadratic planes that need four colours}
\author{Sergi Gal\'an}
\thanks{These results were obtained with the assistance of an AI system (Claude, by Anthropic), which carried
out the searches, the computations and the checks, and helped to draft this note. No mathematician has checked
the note yet. The author is not a professional mathematician and submits it for independent verification.}
\date{Draft, 2 October 2026; remarks added 4 October 2026}
\subjclass[2020]{Primary 05C15, 52C10; Secondary 11R11, 68V15}
\keywords{Hadwiger--Nelson problem, chromatic number of the plane, unit-distance graph, real quadratic field,
SAT solving, DRAT proofs}

\begin{document}

\begin{abstract}
For a subfield $F\subset\R$ let $\chi(F^2)$ be the chromatic number of the graph on $F^2$ joining points at
Euclidean distance $1$. For real quadratic fields $K=\Q(\sqrt d)$ the values known so far were $2$ and $3$, and
a field can need four colours only when $d\equiv 11\pmod{12}$. We show that
$\chi(\Q(\sqrt d)^2)=4$ for $d=11$, 23, 35, 59, 71, 95, 119, 131, 155, 179, 191, 239, 251, 263, 359, 431, 443, 455, 491, 599,
611, 791, 851, 911, 935, 959, and that
$4\le\chi(\Q(\sqrt{47})^2)\le5$. With the known results, $\chi(\Q(\sqrt d)^2)$ is now determined for every
squarefree $d<83$ except $d=47$. Each lower bound is a triangle-free, vertex-critical unit-distance graph with
coordinates in $\Q(\sqrt d)$, with between $71$ and $1\,404$ vertices, that has no proper $3$-colouring: a computer
proof, by a SAT solver whose DRAT proofs are checked by \texttt{drat-trim}, for two encodings written by
separate code. For ten of the fields the whole theorem is also proved in Lean~4 with Mathlib; for the others
(for $d=47$, the lower bound) the Lean proof assumes only that the colouring formula is unsatisfiable, which
the verified checker \texttt{cake\_lpr} confirms. The upper bounds are reductions modulo a prime due to Moorhouse and Fischer. We also explain why
some searches fail: when every direction is integral at certain primes above $2$, $3$ or $5$, reduction there
colours every graph built from those directions with two or three colours. (Added 4 October: a companion draft
now proves $\chi(\Q(\sqrt d)^2)\ge4$ for every $d\equiv11\pmod{12}$ by hand, through characters; the graphs here
remain an explicit and independent proof for these $27$ fields.)
\end{abstract}

\maketitle

\section{Introduction}

The Hadwiger--Nelson problem asks for the chromatic number $\chi(\R^2)$ of the plane; since de Grey
\cite{deGrey} it is known to be $5$, $6$ or $7$. A natural relative restricts the coordinates to a subfield
$F\subset\R$: $\chi(F^2)$ is the least number of colours for the points of $F^2$ such that points at distance
$1$ get different colours. Woodall \cite{Woodall} showed $\chi(\Q^2)=2$. For a real quadratic field
$K=\Q(\sqrt d)$, $d\ge2$ squarefree, the following was known.

\begin{itemize}
\item $\chi(K^2)=2$ if $d\equiv1,2\pmod 4$ (Johnson \cite{Johnson1987}, see \cite{Payne}; Moorhouse
  \cite[Lemma~8.4, Theorem~8.5]{Moorhouse}).
\item $\chi(K^2)\ge3$ if $d\equiv3\pmod 4$ (Fischer \cite[Theorem~8]{Fischer1990}; Moorhouse
  \cite[Theorem~8.6]{Moorhouse} for prime $d$, by an odd cycle from Pell's equation). Directly: writing a
  vector $(x,y)$ as $x+iy$, the unit vector $z=(1+i\sqrt d)/(1-i\sqrt d)$ satisfies
  $\frac{d+1}4(z+\bar z)+\frac{d-1}2\cdot1=0$, a closed walk of odd length $d$.
\item $\chi(K^2)\le3$ if $d\equiv0,1\pmod3$ (Fischer \cite[Theorem~9]{Fischer1990}; Moorhouse
  \cite[Corollary~8.3]{Moorhouse}).
\item $\chi(K^2)\le4$ if $d\equiv0,1,2,4\pmod7$ (Moorhouse \cite[Corollary~8.3]{Moorhouse}) or
  $d\equiv3\pmod8$ (Fischer \cite[Theorem~10]{Fischer1990}; also \cite[\texttt{RESULTS.md}, \S0, item~3]{MM}). Moorhouse's
  Theorem~8.1 states $\chi(K^2)\le4$ unless $d\equiv47,59,83\pmod{84}$.
\end{itemize}

Madore \cite{Madore} proved $\chi(\Q(\sqrt3)^2)=\chi(\Q(\sqrt7)^2)=3$. Together these give $\chi(K^2)=3$ when
$d\equiv3\pmod4$ and $d\not\equiv2\pmod3$; so a real quadratic field can need four colours only if
$d\equiv11\pmod{12}$, and for those the known lower bound was~$3$. As far as we could find, no real quadratic
field was known to need four colours. Moorhouse asked in a 2010 talk \cite{MoorhouseTalk} what can be said
about $\chi(\Q(\sqrt d)^2)$, and about $\chi(\Q(\sqrt{47})^2)$ in particular.

\begin{theorem}\label{thm:main}
$\chi(\Q(\sqrt d)^2)=4$ for
\[\begin{aligned}d={}&11, 23, 35, 59, 71, 95, 119, 131, 155, 179, 191, 239,\\
&251, 263, 359, 431, 443, 455, 491, 599, 611, 791, 851, 911, 935, 959,\end{aligned}\]
and $4\le\chi(\Q(\sqrt{47})^2)\le5$.
\end{theorem}

\begin{corollary}\label{cor:small}
For squarefree $2\le d<83$, $\chi(\Q(\sqrt d)^2)$ is $2$ if $d\equiv1,2\pmod4$, $3$ if $d\equiv3\pmod4$ and
$d\not\equiv2\pmod3$, and $4$ if $d\in\{11,23,35,59,71\}$. The only remaining value is $d=47$, with
$4\le\chi\le5$.
\end{corollary}

\begin{remark}[added 4 October 2026]
Later drafts go further. \cite{FC} proves $\chi(\Q(\sqrt d)^2)\ge4$ for every $d\equiv11\pmod{12}$, by hand, and
\cite{Winding} proves $\chi(\Q(\sqrt d)^2)=4$ for $d=83,107,203$ and $4\le\chi(\Q(\sqrt{143})^2)\le5$ by exact
certificates; so $\chi(\Q(\sqrt d)^2)$ is determined for every squarefree $d<143$ except $d=47$, and equals $4$ for
every $d\equiv11\pmod{12}$ except possibly $d\equiv47,143,167\pmod{168}$. Those proofs go through characters and
do not exhibit graphs; the graphs here remain an independent, explicit proof for the $27$ fields of
Theorem~\ref{thm:main}.
\end{remark}

The upper bounds in Theorem~\ref{thm:main} are known (Section~\ref{sec:upper}). Each lower bound is a finite
graph: a unit-distance graph with coordinates in $\Q(\sqrt d)$ that is not $3$-colourable. That it is not
$3$-colourable is a computer proof (Section~\ref{sec:cert}); the rest of each certificate is checked in exact
integer arithmetic. The graphs, the formulas, the solver logs, the checker and the search code are in the
repository \cite{Repo}.

The graphs are triangle-free: a unit triangle needs $\sqrt{-3}\in K(i)$, which fails for $d\neq3$. As far as
we know, no triangle-free $5$-chromatic unit-distance graph is known (\cite[\texttt{RESULTS.md}, \S0, item~3]{MM}
calls this an open problem);
since $\chi(\Q(\sqrt{47})^2)\le5$ and $\sqrt3\notin\Q(\sqrt{47})$, deciding whether $\chi(\Q(\sqrt{47})^2)$
is $4$ or $5$ touches that question.

\section{The graphs}\label{sec:graphs}

Fix $d$ and a positive integer $D$. A point is written $[a,b,c,e]$ with $a,b,c,e\in\Z$, and stands for
\[
  \Bigl(\frac{a+b\sqrt d}{D},\ \frac{c+e\sqrt d}{D}\Bigr)\in K^2.
\]
Since $1$ and $\sqrt d$ are linearly independent over $\Q$, two points are at distance $1$ exactly when their
differences satisfy the two integer identities
\begin{equation}\label{eq:unit}
  (\Delta a)^2+d\,(\Delta b)^2+(\Delta c)^2+d\,(\Delta e)^2=D^2,\qquad \Delta a\,\Delta b+\Delta c\,\Delta e=0 .
\end{equation}
So every edge is checked exactly. Table~\ref{tab:graphs} lists the graphs, and Figure~\ref{fig:q11} shows the
one over $\Q(\sqrt{11})$.

\begin{table}[ht]
\centering
% written by make_table.py from data/quadratic_planes/
\begin{tabular}{rrrrcr}
\hline
$d$ & $D$ & vertices & edges & degrees & directions used\\
\hline
11 & 30 & 94 & 214 & 3--32 & 68\\
23 & 120 & 660 & 1\,727 & 3--60 & 100\\
35 & 390 & 580 & 1\,501 & 3--107 & 218\\
47 & 240 & 872 & 2\,280 & 3--76 & 102\\
59 & 210 & 406 & 993 & 3--48 & 66\\
71 & 120 & 611 & 1\,557 & 3--85 & 128\\
119 & 240 & 399 & 1\,019 & 3--62 & 106\\
131 & 390 & 356 & 804 & 3--61 & 126\\
179 & 390 & 259 & 622 & 3--39 & 90\\
191 & 240 & 100 & 224 & 3--23 & 52\\
239 & 480 & 355 & 888 & 3--47 & 80\\
359 & 600 & 715 & 1\,851 & 3--63 & 74\\
431 & 600 & 331 & 764 & 3--28 & 66\\
\hline
\end{tabular}

\medskip
\caption{The graphs of Theorem~\ref{thm:main}. $D$ is the common denominator; ``directions used'' counts the
unit vectors $v-u$ over the edges $uv$, with both signs.}\label{tab:graphs}
\end{table}

\begin{figure}[ht]
\centering
% written by make_figure.py from data/quadratic_planes/q11.json
\begin{tikzpicture}[x=3.2cm, y=3.2cm]
\draw[gray!55, densely dashed, line width=0.35pt] (0,0) circle (1);
\draw[gray!65, line width=0.22pt] (-0.90167,-0.18945) -- (-0.06834,0.36332);
\draw[gray!65, line width=0.22pt] (-0.90167,-0.18945) -- (0.09833,-0.18945);
\draw[gray!65, line width=0.22pt] (-0.90167,-0.18945) -- (-0.63333,0.77388);
\draw[gray!65, line width=0.22pt] (-0.90167,-0.18945) -- (-1.71944,0.38610);
\draw[gray!65, line width=0.22pt] (-0.90167,-0.18945) -- (-0.14556,-0.84389);
\draw[gray!65, line width=0.22pt] (-1.71944,0.38610) -- (-1.45110,1.34943);
\draw[gray!65, line width=0.22pt] (-1.71944,0.38610) -- (-0.96332,-0.26834);
\draw[gray!65, line width=0.22pt] (-1.71944,0.38610) -- (-0.75611,0.65444);
\draw[gray!65, line width=0.22pt] (-1.33166,-0.16332) -- (-0.38945,-0.49833);
\draw[gray!65, line width=0.22pt] (-1.33166,-0.16332) -- (-0.57555,-0.81777);
\draw[gray!65, line width=0.22pt] (-1.33166,-0.16332) -- (-0.75611,0.65444);
\draw[gray!65, line width=0.22pt] (-1.00000,0.00000) -- (-0.90000,0.99499);
\draw[gray!65, line width=0.22pt] (-1.00000,0.00000) -- (-0.73166,0.96332);
\draw[gray!65, line width=0.22pt] (-1.00000,0.00000) -- (-0.06834,0.36332);
\draw[gray!65, line width=0.22pt] (-1.00000,0.00000) -- (0.00000,0.00000);
\draw[gray!65, line width=0.22pt] (-1.00000,0.00000) -- (-0.02278,0.21223);
\draw[gray!65, line width=0.22pt] (-0.85611,-0.34055) -- (-0.75611,0.65444);
\draw[gray!65, line width=0.22pt] (-0.85611,-0.34055) -- (-0.02278,0.21223);
\draw[gray!65, line width=0.22pt] (-0.85611,-0.34055) -- (-0.58777,0.62278);
\draw[gray!65, line width=0.22pt] (-0.90000,0.99499) -- (0.07722,1.20721);
\draw[gray!65, line width=0.22pt] (-0.90000,0.99499) -- (0.03166,0.63166);
\draw[gray!65, line width=0.22pt] (-0.90000,0.99499) -- (-0.14389,0.34055);
\draw[gray!65, line width=0.22pt] (-0.97722,-0.21223) -- (-0.04556,0.15110);
\draw[gray!65, line width=0.22pt] (-0.97722,-0.21223) -- (-0.14389,0.34055);
\draw[gray!65, line width=0.22pt] (-0.97722,-0.21223) -- (0.00000,0.00000);
\draw[gray!65, line width=0.22pt] (-0.75611,0.65444) -- (-0.48777,1.61777);
\draw[gray!65, line width=0.22pt] (-0.75611,0.65444) -- (0.07722,1.20721);
\draw[gray!65, line width=0.22pt] (-0.75611,0.65444) -- (0.24389,0.65444);
\draw[gray!65, line width=0.22pt] (-0.75611,0.65444) -- (0.00000,0.00000);
\draw[gray!65, line width=0.22pt] (-0.83333,-0.55277) -- (-1.14222,0.39833);
\draw[gray!65, line width=0.22pt] (-0.83333,-0.55277) -- (-0.25778,0.26500);
\draw[gray!65, line width=0.22pt] (-0.83333,-0.55277) -- (0.09833,-0.18945);
\draw[gray!65, line width=0.22pt] (-0.83333,-0.55277) -- (-0.20000,-1.32665);
\draw[gray!65, line width=0.22pt] (-0.83333,-0.55277) -- (0.00000,0.00000);
\draw[gray!65, line width=0.22pt] (-0.62278,-0.58777) -- (-0.93166,0.36332);
\draw[gray!65, line width=0.22pt] (-0.62278,-0.58777) -- (0.30888,-0.95110);
\draw[gray!65, line width=0.22pt] (-0.62278,-0.58777) -- (-0.02278,0.21223);
\draw[gray!65, line width=0.22pt] (-1.14222,0.39833) -- (-0.30888,0.95110);
\draw[gray!65, line width=0.22pt] (-1.14222,0.39833) -- (-1.45110,1.34943);
\draw[gray!65, line width=0.22pt] (-1.14222,0.39833) -- (-0.50888,-0.37555);
\draw[gray!65, line width=0.22pt] (-1.14222,0.39833) -- (-0.14389,0.34055);
\draw[gray!65, line width=0.22pt] (-0.25778,0.26500) -- (0.31777,1.08276);
\draw[gray!65, line width=0.22pt] (-0.25778,0.26500) -- (0.07722,1.20721);
\draw[gray!65, line width=0.22pt] (-0.25778,0.26500) -- (0.37555,-0.50888);
\draw[gray!65, line width=0.22pt] (-0.25778,0.26500) -- (0.57555,0.81777);
\draw[gray!65, line width=0.22pt] (-0.63333,0.77388) -- (0.20000,1.32665);
\draw[gray!65, line width=0.22pt] (-0.63333,0.77388) -- (-1.45110,1.34943);
\draw[gray!65, line width=0.22pt] (-0.63333,0.77388) -- (0.00000,0.00000);
\draw[gray!65, line width=0.22pt] (-0.63333,0.77388) -- (0.31777,1.08276);
\draw[gray!65, line width=0.22pt] (-0.63333,0.77388) -- (0.36667,0.77388);
\draw[gray!65, line width=0.22pt] (-0.93166,0.36332) -- (0.03166,0.63166);
\draw[gray!65, line width=0.22pt] (-0.93166,0.36332) -- (0.00000,0.00000);
\draw[gray!65, line width=0.22pt] (-0.93166,0.36332) -- (0.06834,0.36332);
\draw[gray!65, line width=0.22pt] (-0.93166,0.36332) -- (-1.29499,1.29499);
\draw[gray!65, line width=0.22pt] (-0.93166,-0.36332) -- (0.03166,-0.63166);
\draw[gray!65, line width=0.22pt] (-0.93166,-0.36332) -- (0.04556,-0.15110);
\draw[gray!65, line width=0.22pt] (-0.93166,-0.36332) -- (0.00000,0.00000);
\draw[gray!65, line width=0.22pt] (-1.45110,1.34943) -- (-0.48777,1.61777);
\draw[gray!65, line width=0.22pt] (-1.45110,1.34943) -- (-0.81777,0.57555);
\draw[gray!65, line width=0.22pt] (0.31777,1.08276) -- (1.29499,1.29499);
\draw[gray!65, line width=0.22pt] (0.31777,1.08276) -- (0.95110,0.30888);
\draw[gray!65, line width=0.22pt] (0.31777,1.08276) -- (-0.04556,0.15110);
\draw[gray!65, line width=0.22pt] (-0.38945,-0.49833) -- (0.55277,-0.83333);
\draw[gray!65, line width=0.22pt] (-0.38945,-0.49833) -- (0.36667,-1.15277);
\draw[gray!65, line width=0.22pt] (-0.02445,1.61777) -- (0.24389,0.65444);
\draw[gray!65, line width=0.22pt] (-0.02445,1.61777) -- (0.57555,0.81777);
\draw[gray!65, line width=0.22pt] (-0.02445,1.61777) -- (0.73166,0.96332);
\draw[gray!65, line width=0.22pt] (-0.73166,0.96332) -- (0.06834,0.36332);
\draw[gray!65, line width=0.22pt] (-0.73166,0.96332) -- (0.20000,1.32665);
\draw[gray!65, line width=0.22pt] (-0.73166,0.96332) -- (0.26834,0.96332);
\draw[gray!65, line width=0.22pt] (-0.06834,0.36332) -- (0.93166,0.36332);
\draw[gray!65, line width=0.22pt] (-0.06834,0.36332) -- (0.20000,1.32665);
\draw[gray!65, line width=0.22pt] (-0.06834,0.36332) -- (0.03166,-0.63166);
\draw[gray!65, line width=0.22pt] (-0.06834,0.36332) -- (0.73166,0.96332);
\draw[gray!65, line width=0.22pt] (-0.58777,0.62278) -- (0.04556,-0.15110);
\draw[gray!65, line width=0.22pt] (-0.58777,0.62278) -- (0.36332,0.93166);
\draw[gray!65, line width=0.22pt] (-0.58777,0.62278) -- (-0.48777,1.61777);
\draw[gray!65, line width=0.22pt] (-0.58777,0.62278) -- (-0.95110,-0.30888);
\draw[gray!65, line width=0.22pt] (-0.58777,0.62278) -- (0.41055,0.56500);
\draw[gray!65, line width=0.22pt] (-0.14556,-0.84389) -- (0.18945,0.09833);
\draw[gray!65, line width=0.22pt] (-0.14556,-0.84389) -- (-0.04556,0.15110);
\draw[gray!65, line width=0.22pt] (-0.14556,-0.84389) -- (-0.96332,-0.26834);
\draw[gray!65, line width=0.22pt] (-1.29499,1.29499) -- (-1.08276,0.31777);
\draw[gray!65, line width=0.22pt] (-1.29499,1.29499) -- (-0.36332,0.93166);
\draw[gray!65, line width=0.22pt] (-0.96332,-0.26834) -- (-0.03166,-0.63166);
\draw[gray!65, line width=0.22pt] (-0.96332,-0.26834) -- (-1.32665,-1.20000);
\draw[gray!65, line width=0.22pt] (-0.96332,-0.26834) -- (-0.62832,0.67388);
\draw[gray!65, line width=0.22pt] (-0.96332,-0.26834) -- (0.00000,0.00000);
\draw[gray!65, line width=0.22pt] (-0.96332,-0.26834) -- (-0.96332,0.73166);
\draw[gray!65, line width=0.22pt] (-0.96332,-0.73166) -- (-1.32665,0.20000);
\draw[gray!65, line width=0.22pt] (-0.96332,-0.73166) -- (0.00000,-1.00000);
\draw[gray!65, line width=0.22pt] (-0.96332,-0.73166) -- (0.03166,-0.63166);
\draw[gray!65, line width=0.22pt] (-0.96332,-0.73166) -- (-0.96332,0.26834);
\draw[gray!65, line width=0.22pt] (-0.96332,0.73166) -- (-0.03166,0.36834);
\draw[gray!65, line width=0.22pt] (-0.96332,0.73166) -- (-0.36332,-0.06834);
\draw[gray!65, line width=0.22pt] (-0.96332,0.73166) -- (0.03166,0.63166);
\draw[gray!65, line width=0.22pt] (-0.96332,0.26834) -- (-0.03166,-0.09499);
\draw[gray!65, line width=0.22pt] (-0.96332,0.26834) -- (0.00000,0.00000);
\draw[gray!65, line width=0.22pt] (0.03166,-0.63166) -- (0.99499,-0.90000);
\draw[gray!65, line width=0.22pt] (0.03166,-0.63166) -- (0.96332,-0.26834);
\draw[gray!65, line width=0.22pt] (0.03166,0.63166) -- (0.96332,0.26834);
\draw[gray!65, line width=0.22pt] (0.36332,0.93166) -- (0.46332,1.92665);
\draw[gray!65, line width=0.22pt] (0.36332,0.93166) -- (1.29499,1.29499);
\draw[gray!65, line width=0.22pt] (0.36332,0.93166) -- (0.00000,0.00000);
\draw[gray!65, line width=0.22pt] (-0.48777,1.61777) -- (0.14556,0.84389);
\draw[gray!65, line width=0.22pt] (-0.48777,1.61777) -- (0.46332,1.92665);
\draw[gray!65, line width=0.22pt] (-0.48777,1.61777) -- (0.26834,0.96332);
\draw[gray!65, line width=0.22pt] (-0.04556,0.15110) -- (0.93166,0.36332);
\draw[gray!65, line width=0.22pt] (-0.04556,0.15110) -- (0.58777,-0.62278);
\draw[gray!65, line width=0.22pt] (0.09833,-0.18945) -- (0.36667,-1.15277);
\draw[gray!65, line width=0.22pt] (0.09833,-0.18945) -- (0.93166,0.36332);
\draw[gray!65, line width=0.22pt] (0.09833,-0.18945) -- (0.36667,0.77388);
\draw[gray!65, line width=0.22pt] (-0.20000,-1.32665) -- (-0.50888,-0.37555);
\draw[gray!65, line width=0.22pt] (-0.20000,-1.32665) -- (0.37555,-0.50888);
\draw[gray!65, line width=0.22pt] (0.46332,1.92665) -- (0.73166,0.96332);
\draw[gray!65, line width=0.22pt] (-0.68610,-0.85110) -- (-0.03166,-0.09499);
\draw[gray!65, line width=0.22pt] (-0.68610,-0.85110) -- (0.30888,-0.95110);
\draw[gray!65, line width=0.22pt] (-0.68610,-0.85110) -- (-0.99499,0.10000);
\draw[gray!65, line width=0.22pt] (-0.57555,-0.81777) -- (0.37555,-0.50888);
\draw[gray!65, line width=0.22pt] (-0.57555,-0.81777) -- (0.25778,-0.26500);
\draw[gray!65, line width=0.22pt] (-0.57555,-0.81777) -- (0.00000,0.00000);
\draw[gray!65, line width=0.22pt] (-0.57555,-0.81777) -- (0.36667,-1.15277);
\draw[gray!65, line width=0.22pt] (-0.02278,0.21223) -- (0.97722,0.21223);
\draw[gray!65, line width=0.22pt] (-0.02278,0.21223) -- (0.07722,1.20721);
\draw[gray!65, line width=0.22pt] (0.30888,-0.95110) -- (-0.50888,-0.37555);
\draw[gray!65, line width=0.22pt] (0.30888,-0.95110) -- (0.00000,0.00000);
\draw[gray!65, line width=0.22pt] (-0.95110,-0.30888) -- (-0.39833,-1.14222);
\draw[gray!65, line width=0.22pt] (-0.95110,-0.30888) -- (-0.37555,0.50888);
\draw[gray!65, line width=0.22pt] (-0.95110,-0.30888) -- (0.00000,0.00000);
\draw[gray!65, line width=0.22pt] (-0.50888,-0.37555) -- (0.04389,-1.20888);
\draw[gray!65, line width=0.22pt] (-0.50888,-0.37555) -- (-1.32665,0.20000);
\draw[gray!65, line width=0.22pt] (-0.50888,-0.37555) -- (-0.81777,0.57555);
\draw[gray!65, line width=0.22pt] (-0.39833,-1.14222) -- (0.55277,-0.83333);
\draw[gray!65, line width=0.22pt] (-0.39833,-1.14222) -- (0.17722,-0.32445);
\draw[gray!65, line width=0.22pt] (-0.39833,-1.14222) -- (0.37555,-0.50888);
\draw[gray!65, line width=0.22pt] (0.04389,-1.20888) -- (0.99499,-0.90000);
\draw[gray!65, line width=0.22pt] (0.04389,-1.20888) -- (-0.77388,-0.63333);
\draw[gray!65, line width=0.22pt] (0.04389,-1.20888) -- (-0.26500,-0.25778);
\draw[gray!65, line width=0.22pt] (0.37555,-0.50888) -- (0.95110,0.30888);
\draw[gray!65, line width=0.22pt] (-0.14389,0.34055) -- (0.83333,0.55277);
\draw[gray!65, line width=0.22pt] (0.07722,1.20721) -- (0.83333,0.55277);
\draw[gray!65, line width=0.22pt] (-1.32665,-1.20000) -- (-0.36332,-0.93166);
\draw[gray!65, line width=0.22pt] (-1.32665,-1.20000) -- (-0.44222,-0.73333);
\draw[gray!65, line width=0.22pt] (-1.32665,0.20000) -- (-0.77388,-0.63333);
\draw[gray!65, line width=0.22pt] (-1.32665,0.20000) -- (-0.36332,-0.06834);
\draw[gray!65, line width=0.22pt] (-1.32665,0.20000) -- (-0.55277,0.83333);
\draw[gray!65, line width=0.22pt] (-1.32665,0.20000) -- (-0.37555,0.50888);
\draw[gray!65, line width=0.22pt] (-0.99499,0.10000) -- (-0.03166,0.36834);
\draw[gray!65, line width=0.22pt] (-0.99499,0.10000) -- (-0.44222,-0.73333);
\draw[gray!65, line width=0.22pt] (-0.99499,0.10000) -- (0.00000,0.00000);
\draw[gray!65, line width=0.22pt] (-0.77388,-0.63333) -- (0.17722,-0.32445);
\draw[gray!65, line width=0.22pt] (-0.77388,-0.63333) -- (-1.08276,0.31777);
\draw[gray!65, line width=0.22pt] (-0.77388,-0.63333) -- (0.00000,0.00000);
\draw[gray!65, line width=0.22pt] (-0.55277,0.83333) -- (0.41055,0.56500);
\draw[gray!65, line width=0.22pt] (-0.55277,0.83333) -- (0.00000,0.00000);
\draw[gray!65, line width=0.22pt] (-0.44222,-0.73333) -- (0.55277,-0.83333);
\draw[gray!65, line width=0.22pt] (0.00000,-1.00000) -- (0.99499,-0.90000);
\draw[gray!65, line width=0.22pt] (0.00000,-1.00000) -- (-0.36332,-0.06834);
\draw[gray!65, line width=0.22pt] (0.00000,-1.00000) -- (0.00000,0.00000);
\draw[gray!65, line width=0.22pt] (0.00000,0.00000) -- (0.26834,0.96332);
\draw[gray!65, line width=0.22pt] (0.00000,0.00000) -- (0.83333,0.55277);
\draw[gray!65, line width=0.22pt] (0.00000,0.00000) -- (0.93166,0.36332);
\draw[gray!65, line width=0.22pt] (0.00000,0.00000) -- (0.93166,-0.36332);
\draw[gray!65, line width=0.22pt] (0.00000,0.00000) -- (0.97722,0.21223);
\draw[gray!65, line width=0.22pt] (0.00000,0.00000) -- (-0.81777,0.57555);
\draw[gray!65, line width=0.22pt] (0.00000,0.00000) -- (0.95110,0.30888);
\draw[gray!65, line width=0.22pt] (0.00000,0.00000) -- (-0.36332,-0.93166);
\draw[gray!65, line width=0.22pt] (0.00000,0.00000) -- (-0.36332,0.93166);
\draw[gray!65, line width=0.22pt] (0.00000,0.00000) -- (0.96332,-0.26834);
\draw[gray!65, line width=0.22pt] (0.00000,0.00000) -- (1.00000,0.00000);
\draw[gray!65, line width=0.22pt] (0.00000,0.00000) -- (0.96332,0.26834);
\draw[gray!65, line width=0.22pt] (0.00000,0.00000) -- (-0.30888,0.95110);
\draw[gray!65, line width=0.22pt] (0.00000,0.00000) -- (0.57555,0.81777);
\draw[gray!65, line width=0.22pt] (0.00000,0.00000) -- (0.55277,-0.83333);
\draw[gray!65, line width=0.22pt] (0.55277,-0.83333) -- (-0.26500,-0.25778);
\draw[gray!65, line width=0.22pt] (0.55277,-0.83333) -- (0.18945,0.09833);
\draw[gray!65, line width=0.22pt] (0.99499,-0.90000) -- (0.17722,-0.32445);
\draw[gray!65, line width=0.22pt] (-0.81777,0.57555) -- (-0.26500,-0.25778);
\draw[gray!65, line width=0.22pt] (-0.81777,0.57555) -- (0.14556,0.84389);
\draw[gray!65, line width=0.22pt] (-0.37555,0.50888) -- (0.17722,-0.32445);
\draw[gray!65, line width=0.22pt] (-0.37555,0.50888) -- (0.20000,1.32665);
\draw[gray!65, line width=0.22pt] (-0.37555,0.50888) -- (0.25778,-0.26500);
\draw[gray!65, line width=0.22pt] (-0.37555,0.50888) -- (0.57555,0.81777);
\draw[gray!65, line width=0.22pt] (-0.26500,-0.25778) -- (-0.62832,0.67388);
\draw[gray!65, line width=0.22pt] (-0.26500,-0.25778) -- (-1.08276,0.31777);
\draw[gray!65, line width=0.22pt] (0.17722,-0.32445) -- (0.95110,0.30888);
\draw[gray!65, line width=0.22pt] (0.95110,0.30888) -- (0.58777,-0.62278);
\draw[gray!65, line width=0.22pt] (-1.08276,0.31777) -- (-0.30888,0.95110);
\draw[gray!65, line width=0.22pt] (0.24389,0.65444) -- (1.00000,0.00000);
\draw[gray!65, line width=0.22pt] (0.20000,1.32665) -- (1.20000,1.32665);
\draw[gray!65, line width=0.22pt] (0.20000,1.32665) -- (0.83333,0.55277);
\draw[gray!65, line width=0.22pt] (0.04556,-0.15110) -- (0.14556,0.84389);
\draw[gray!65, line width=0.22pt] (0.04556,-0.15110) -- (0.97722,0.21223);
\draw[gray!65, line width=0.22pt] (-0.36332,-0.93166) -- (0.58777,-0.62278);
\draw[gray!65, line width=0.22pt] (-0.36332,-0.06834) -- (-0.36332,0.93166);
\draw[gray!65, line width=0.22pt] (-0.36332,-0.06834) -- (0.41055,0.56500);
\draw[gray!65, line width=0.22pt] (-0.36332,0.93166) -- (0.18945,0.09833);
\draw[gray!65, line width=0.22pt] (-0.03166,-0.63166) -- (-0.03166,0.36834);
\draw[gray!65, line width=0.22pt] (-0.03166,-0.63166) -- (0.06834,0.36332);
\draw[gray!65, line width=0.22pt] (-0.03166,-0.63166) -- (0.93166,-0.36332);
\draw[gray!65, line width=0.22pt] (-0.03166,0.36834) -- (0.96332,0.26834);
\draw[gray!65, line width=0.22pt] (-0.03166,-0.09499) -- (0.93166,-0.36332);
\draw[gray!65, line width=0.22pt] (0.18945,0.09833) -- (-0.62832,0.67388);
\draw[gray!65, line width=0.22pt] (0.41055,0.56500) -- (0.96332,-0.26834);
\draw[gray!65, line width=0.22pt] (0.96332,0.26834) -- (0.14556,0.84389);
\draw[gray!65, line width=0.22pt] (1.29499,1.29499) -- (0.93166,0.36332);
\draw[gray!65, line width=0.22pt] (-0.62832,0.67388) -- (0.36667,0.77388);
\draw[gray!65, line width=0.22pt] (0.36667,0.77388) -- (1.00000,0.00000);
\draw[gray!65, line width=0.22pt] (0.36667,0.77388) -- (1.20000,1.32665);
\draw[gray!65, line width=0.22pt] (0.06834,0.36332) -- (1.00000,0.00000);
\draw[gray!65, line width=0.22pt] (0.73166,0.96332) -- (1.00000,0.00000);
\draw[gray!65, line width=0.22pt] (0.26834,0.96332) -- (1.20000,1.32665);
\draw[gray!65, line width=0.22pt] (0.93166,0.36332) -- (1.20000,1.32665);
\draw[gray!65, line width=0.22pt] (0.25778,-0.26500) -- (0.83333,0.55277);
\node[draw=black, line width=0.35pt, fill=white, inner sep=0pt, diamond, minimum size=4.6pt] at (-0.90167,-0.18945) {};
\node[draw=black, line width=0.35pt, fill=black, inner sep=0pt, regular polygon, regular polygon sides=3, minimum size=5pt] at (-1.71944,0.38610) {};
\node[draw=black, line width=0.35pt, fill=white, inner sep=0pt, diamond, minimum size=4.6pt] at (-1.33166,-0.16332) {};
\node[draw=black, line width=0.35pt, fill=white, inner sep=0pt, diamond, minimum size=4.6pt] at (-1.00000,0.00000) {};
\node[draw=black, line width=0.35pt, fill=white, inner sep=0pt, diamond, minimum size=4.6pt] at (-0.85611,-0.34055) {};
\node[draw=black, line width=0.35pt, fill=black, inner sep=0pt, regular polygon, regular polygon sides=3, minimum size=5pt] at (-0.90000,0.99499) {};
\node[draw=black, line width=0.35pt, fill=white, inner sep=0pt, diamond, minimum size=4.6pt] at (-0.97722,-0.21223) {};
\node[draw=black, line width=0.35pt, fill=white, inner sep=0pt, rectangle, minimum size=3.3pt] at (-0.75611,0.65444) {};
\node[draw=black, line width=0.35pt, fill=white, inner sep=0pt, diamond, minimum size=4.6pt] at (-0.83333,-0.55277) {};
\node[draw=black, line width=0.35pt, fill=white, inner sep=0pt, diamond, minimum size=4.6pt] at (-0.62278,-0.58777) {};
\node[draw=black, line width=0.35pt, fill=black, inner sep=0pt, regular polygon, regular polygon sides=3, minimum size=5pt] at (-1.14222,0.39833) {};
\node[draw=black, line width=0.35pt, fill=black, inner sep=0pt, regular polygon, regular polygon sides=3, minimum size=5pt] at (-0.25778,0.26500) {};
\node[draw=black, line width=0.35pt, fill=black, inner sep=0pt, regular polygon, regular polygon sides=3, minimum size=5pt] at (-0.63333,0.77388) {};
\node[draw=black, line width=0.35pt, fill=black, inner sep=0pt, regular polygon, regular polygon sides=3, minimum size=5pt] at (-0.93166,0.36332) {};
\node[draw=black, line width=0.35pt, fill=white, inner sep=0pt, diamond, minimum size=4.6pt] at (-0.93166,-0.36332) {};
\node[draw=black, line width=0.35pt, fill=white, inner sep=0pt, diamond, minimum size=4.6pt] at (-1.45110,1.34943) {};
\node[draw=black, line width=0.35pt, fill=white, inner sep=0pt, diamond, minimum size=4.6pt] at (0.31777,1.08276) {};
\node[draw=black, line width=0.35pt, fill=black, inner sep=0pt, regular polygon, regular polygon sides=3, minimum size=5pt] at (-0.38945,-0.49833) {};
\node[draw=black, line width=0.35pt, fill=white, inner sep=0pt, diamond, minimum size=4.6pt] at (-0.02445,1.61777) {};
\node[draw=black, line width=0.35pt, fill=black, inner sep=0pt, regular polygon, regular polygon sides=3, minimum size=5pt] at (-0.73166,0.96332) {};
\node[draw=black, line width=0.35pt, fill=black, inner sep=0pt, regular polygon, regular polygon sides=3, minimum size=5pt] at (-0.06834,0.36332) {};
\node[draw=black, line width=0.35pt, fill=black, inner sep=0pt, regular polygon, regular polygon sides=3, minimum size=5pt] at (-0.58777,0.62278) {};
\node[draw=black, line width=0.35pt, fill=black, inner sep=0pt, regular polygon, regular polygon sides=3, minimum size=5pt] at (-0.14556,-0.84389) {};
\node[draw=black, line width=0.35pt, fill=white, inner sep=0pt, diamond, minimum size=4.6pt] at (-1.29499,1.29499) {};
\node[draw=black, line width=0.35pt, fill=white, inner sep=0pt, diamond, minimum size=4.6pt] at (-0.96332,-0.26834) {};
\node[draw=black, line width=0.35pt, fill=white, inner sep=0pt, diamond, minimum size=4.6pt] at (-0.96332,-0.73166) {};
\node[draw=black, line width=0.35pt, fill=black, inner sep=0pt, regular polygon, regular polygon sides=3, minimum size=5pt] at (-0.96332,0.73166) {};
\node[draw=black, line width=0.35pt, fill=black, inner sep=0pt, regular polygon, regular polygon sides=3, minimum size=5pt] at (-0.96332,0.26834) {};
\node[draw=black, line width=0.35pt, fill=white, inner sep=0pt, rectangle, minimum size=3.3pt] at (0.03166,-0.63166) {};
\node[draw=black, line width=0.35pt, fill=white, inner sep=0pt, diamond, minimum size=4.6pt] at (0.03166,0.63166) {};
\node[draw=black, line width=0.35pt, fill=white, inner sep=0pt, diamond, minimum size=4.6pt] at (0.36332,0.93166) {};
\node[draw=black, line width=0.35pt, fill=black, inner sep=0pt, circle, minimum size=3.6pt] at (-0.48777,1.61777) {};
\node[draw=black, line width=0.35pt, fill=white, inner sep=0pt, rectangle, minimum size=3.3pt] at (-0.04556,0.15110) {};
\node[draw=black, line width=0.35pt, fill=black, inner sep=0pt, regular polygon, regular polygon sides=3, minimum size=5pt] at (0.09833,-0.18945) {};
\node[draw=black, line width=0.35pt, fill=black, inner sep=0pt, regular polygon, regular polygon sides=3, minimum size=5pt] at (-0.20000,-1.32665) {};
\node[draw=black, line width=0.35pt, fill=black, inner sep=0pt, regular polygon, regular polygon sides=3, minimum size=5pt] at (0.46332,1.92665) {};
\node[draw=black, line width=0.35pt, fill=white, inner sep=0pt, diamond, minimum size=4.6pt] at (-0.68610,-0.85110) {};
\node[draw=black, line width=0.35pt, fill=black, inner sep=0pt, regular polygon, regular polygon sides=3, minimum size=5pt] at (-0.57555,-0.81777) {};
\node[draw=black, line width=0.35pt, fill=black, inner sep=0pt, regular polygon, regular polygon sides=3, minimum size=5pt] at (-0.02278,0.21223) {};
\node[draw=black, line width=0.35pt, fill=black, inner sep=0pt, regular polygon, regular polygon sides=3, minimum size=5pt] at (0.30888,-0.95110) {};
\node[draw=black, line width=0.35pt, fill=white, inner sep=0pt, diamond, minimum size=4.6pt] at (-0.95110,-0.30888) {};
\node[draw=black, line width=0.35pt, fill=white, inner sep=0pt, diamond, minimum size=4.6pt] at (-0.50888,-0.37555) {};
\node[draw=black, line width=0.35pt, fill=black, inner sep=0pt, regular polygon, regular polygon sides=3, minimum size=5pt] at (-0.39833,-1.14222) {};
\node[draw=black, line width=0.35pt, fill=black, inner sep=0pt, regular polygon, regular polygon sides=3, minimum size=5pt] at (0.04389,-1.20888) {};
\node[draw=black, line width=0.35pt, fill=white, inner sep=0pt, diamond, minimum size=4.6pt] at (0.37555,-0.50888) {};
\node[draw=black, line width=0.35pt, fill=white, inner sep=0pt, rectangle, minimum size=3.3pt] at (-0.14389,0.34055) {};
\node[draw=black, line width=0.35pt, fill=white, inner sep=0pt, diamond, minimum size=4.6pt] at (0.07722,1.20721) {};
\node[draw=black, line width=0.35pt, fill=black, inner sep=0pt, regular polygon, regular polygon sides=3, minimum size=5pt] at (-1.32665,-1.20000) {};
\node[draw=black, line width=0.35pt, fill=black, inner sep=0pt, regular polygon, regular polygon sides=3, minimum size=5pt] at (-1.32665,0.20000) {};
\node[draw=black, line width=0.35pt, fill=black, inner sep=0pt, regular polygon, regular polygon sides=3, minimum size=5pt] at (-0.99499,0.10000) {};
\node[draw=black, line width=0.35pt, fill=white, inner sep=0pt, diamond, minimum size=4.6pt] at (-0.77388,-0.63333) {};
\node[draw=black, line width=0.35pt, fill=white, inner sep=0pt, diamond, minimum size=4.6pt] at (-0.55277,0.83333) {};
\node[draw=black, line width=0.35pt, fill=white, inner sep=0pt, diamond, minimum size=4.6pt] at (-0.44222,-0.73333) {};
\node[draw=black, line width=0.35pt, fill=black, inner sep=0pt, regular polygon, regular polygon sides=3, minimum size=5pt] at (0.00000,-1.00000) {};
\node[draw=black, line width=0.35pt, fill=black, inner sep=0pt, circle, minimum size=3.6pt] at (0.00000,0.00000) {};
\node[draw=black, line width=0.35pt, fill=white, inner sep=0pt, rectangle, minimum size=3.3pt] at (0.55277,-0.83333) {};
\node[draw=black, line width=0.35pt, fill=white, inner sep=0pt, diamond, minimum size=4.6pt] at (0.99499,-0.90000) {};
\node[draw=black, line width=0.35pt, fill=black, inner sep=0pt, regular polygon, regular polygon sides=3, minimum size=5pt] at (-0.81777,0.57555) {};
\node[draw=black, line width=0.35pt, fill=black, inner sep=0pt, circle, minimum size=3.6pt] at (-0.37555,0.50888) {};
\node[draw=black, line width=0.35pt, fill=white, inner sep=0pt, diamond, minimum size=4.6pt] at (-0.26500,-0.25778) {};
\node[draw=black, line width=0.35pt, fill=white, inner sep=0pt, rectangle, minimum size=3.3pt] at (0.17722,-0.32445) {};
\node[draw=black, line width=0.35pt, fill=black, inner sep=0pt, regular polygon, regular polygon sides=3, minimum size=5pt] at (0.95110,0.30888) {};
\node[draw=black, line width=0.35pt, fill=black, inner sep=0pt, regular polygon, regular polygon sides=3, minimum size=5pt] at (-1.08276,0.31777) {};
\node[draw=black, line width=0.35pt, fill=white, inner sep=0pt, diamond, minimum size=4.6pt] at (-0.30888,0.95110) {};
\node[draw=black, line width=0.35pt, fill=black, inner sep=0pt, regular polygon, regular polygon sides=3, minimum size=5pt] at (0.24389,0.65444) {};
\node[draw=black, line width=0.35pt, fill=white, inner sep=0pt, rectangle, minimum size=3.3pt] at (0.57555,0.81777) {};
\node[draw=black, line width=0.35pt, fill=white, inner sep=0pt, diamond, minimum size=4.6pt] at (0.20000,1.32665) {};
\node[draw=black, line width=0.35pt, fill=white, inner sep=0pt, rectangle, minimum size=3.3pt] at (0.04556,-0.15110) {};
\node[draw=black, line width=0.35pt, fill=white, inner sep=0pt, diamond, minimum size=4.6pt] at (-0.36332,-0.93166) {};
\node[draw=black, line width=0.35pt, fill=white, inner sep=0pt, diamond, minimum size=4.6pt] at (-0.36332,-0.06834) {};
\node[draw=black, line width=0.35pt, fill=black, inner sep=0pt, regular polygon, regular polygon sides=3, minimum size=5pt] at (-0.36332,0.93166) {};
\node[draw=black, line width=0.35pt, fill=black, inner sep=0pt, regular polygon, regular polygon sides=3, minimum size=5pt] at (-0.03166,-0.63166) {};
\node[draw=black, line width=0.35pt, fill=white, inner sep=0pt, diamond, minimum size=4.6pt] at (-0.03166,0.36834) {};
\node[draw=black, line width=0.35pt, fill=white, inner sep=0pt, rectangle, minimum size=3.3pt] at (-0.03166,-0.09499) {};
\node[draw=black, line width=0.35pt, fill=white, inner sep=0pt, diamond, minimum size=4.6pt] at (0.18945,0.09833) {};
\node[draw=black, line width=0.35pt, fill=white, inner sep=0pt, rectangle, minimum size=3.3pt] at (0.41055,0.56500) {};
\node[draw=black, line width=0.35pt, fill=white, inner sep=0pt, diamond, minimum size=4.6pt] at (0.96332,-0.26834) {};
\node[draw=black, line width=0.35pt, fill=black, inner sep=0pt, regular polygon, regular polygon sides=3, minimum size=5pt] at (0.96332,0.26834) {};
\node[draw=black, line width=0.35pt, fill=black, inner sep=0pt, regular polygon, regular polygon sides=3, minimum size=5pt] at (1.29499,1.29499) {};
\node[draw=black, line width=0.35pt, fill=black, inner sep=0pt, regular polygon, regular polygon sides=3, minimum size=5pt] at (-0.62832,0.67388) {};
\node[draw=black, line width=0.35pt, fill=white, inner sep=0pt, diamond, minimum size=4.6pt] at (0.14556,0.84389) {};
\node[draw=black, line width=0.35pt, fill=white, inner sep=0pt, diamond, minimum size=4.6pt] at (0.36667,-1.15277) {};
\node[draw=black, line width=0.35pt, fill=white, inner sep=0pt, diamond, minimum size=4.6pt] at (0.36667,0.77388) {};
\node[draw=black, line width=0.35pt, fill=black, inner sep=0pt, circle, minimum size=3.6pt] at (0.58777,-0.62278) {};
\node[draw=black, line width=0.35pt, fill=white, inner sep=0pt, diamond, minimum size=4.6pt] at (0.06834,0.36332) {};
\node[draw=black, line width=0.35pt, fill=black, inner sep=0pt, circle, minimum size=3.6pt] at (0.73166,0.96332) {};
\node[draw=black, line width=0.35pt, fill=white, inner sep=0pt, diamond, minimum size=4.6pt] at (0.26834,0.96332) {};
\node[draw=black, line width=0.35pt, fill=white, inner sep=0pt, diamond, minimum size=4.6pt] at (0.93166,0.36332) {};
\node[draw=black, line width=0.35pt, fill=white, inner sep=0pt, diamond, minimum size=4.6pt] at (0.93166,-0.36332) {};
\node[draw=black, line width=0.35pt, fill=white, inner sep=0pt, diamond, minimum size=4.6pt] at (0.25778,-0.26500) {};
\node[draw=black, line width=0.35pt, fill=black, inner sep=0pt, regular polygon, regular polygon sides=3, minimum size=5pt] at (0.83333,0.55277) {};
\node[draw=black, line width=0.35pt, fill=white, inner sep=0pt, diamond, minimum size=4.6pt] at (0.97722,0.21223) {};
\node[draw=black, line width=0.35pt, fill=white, inner sep=0pt, rectangle, minimum size=3.3pt] at (1.00000,0.00000) {};
\node[draw=black, line width=0.35pt, fill=black, inner sep=0pt, regular polygon, regular polygon sides=3, minimum size=5pt] at (1.20000,1.32665) {};
\end{tikzpicture}

\caption{The graph over $\Q(\sqrt{11})$: $76$ points, joined by the $172$ segments of length exactly $1$ (drawn from
the coordinates rounded to five decimals; \texttt{make\_figure.py} writes the drawing from the data). The marker
shapes show a proper $4$-colouring; there is no proper $3$-colouring. The vertex at the origin has $26$ neighbours, on
the dashed unit circle.}\label{fig:q11}
\end{figure}

For each graph $G$:
\begin{itemize}
\item the edges are all the pairs of its points that satisfy \eqref{eq:unit};
\item there is no triangle, and there is a $4$-cycle, so the girth is $4$;
\item there is a proper $4$-colouring, and no proper $3$-colouring;
\item $G$ is vertex-critical: for every vertex $v$ there is a proper $3$-colouring of $G-v$.
\end{itemize}

\section{The certificates}\label{sec:cert}

For a graph $G$ with vertices $0,\dots,n-1$ and a fixed edge $u_0v_0$, let $\Phi(G)$ be the formula with
variables $x_{v,c}$ ($v$ a vertex, $c\in\{0,1,2\}$), the clauses $x_{v,0}\lor x_{v,1}\lor x_{v,2}$ for each
vertex, $\lnot x_{u,c}\lor\lnot x_{v,c}$ for each edge $uv$ and colour $c$, and the unit clauses $x_{u_0,0}$
and $x_{v_0,1}$. Renaming colours, $G$ is $3$-colourable if and only if $\Phi(G)$ is satisfiable.

For each $d$ the repository stores $\Phi(G)$, with $x_{v,c}$ numbered $3v+c+1$. The solver kissat~4.0.4
answered UNSATISFIABLE and wrote a DRAT proof, which \texttt{drat-trim} verified. The same was done for a second
encoding, written by separate code, with $x_{v,c}$ numbered $cn+v+1$. The DRAT proofs, in kissat's binary format,
have from $1.4$~kB to $17$~MB, and each check takes at most $18$ seconds.

The checker \texttt{scripts/verify\_quadratic\_planes.py}, which uses only the Python standard library, checks
again, for each $d$: that the points are distinct, that every listed edge satisfies \eqref{eq:unit} and every
pair satisfying \eqref{eq:unit} is listed, that there is no triangle and that there is a $4$-cycle, the
stored $4$-colouring, the stored
$3$-colourings of $G-v$ for every $v$, that the stored formula is exactly $\Phi(G)$ (and prints its SHA-256
hash, so that a log names the files it checked), and the colourings of the
finite planes behind the upper bounds. With a path to kissat and \texttt{drat-trim} it also writes the second
encoding, solves it and checks the proof. With a path to \texttt{cake\_lpr} as well, \texttt{drat-trim} writes
each proof in LRAT form and \texttt{cake\_lpr}, a proof checker verified in HOL4 and compiled by the verified
CakeML compiler \cite{cakelpr}, checks it, also for the stored formula; it accepted the proofs for all
twenty-seven graphs. What has to be trusted
is a proof checker (\texttt{drat-trim}, or \texttt{cake\_lpr} and its formal verification), the short code that
writes $\Phi(G)$ from the list of edges, and exact integer arithmetic.

For $d=11,119,131,179,191,251,431,455,911,935$ the whole theorem is also proved in Lean~4 \cite{Lean} with
Mathlib \cite{mathlib} (\texttt{lean/Sqrt$d$.lean} in \cite{Repo}). For each of these fields the kernel checks
the unit distances of the graph, an LRAT proof that $\Phi(G)$ is unsatisfiable, through Mathlib's
\texttt{lrat\_proof}, and the upper bound of Section~\ref{sec:upper}: at $p=7$, from a valuation ring of
$\Q(\sqrt d)$ with residue field $\F_7$ (when $d$ is a nonzero square modulo~$7$, and when $7\mid d$), and for
$d=131,251$ at the place above~$2$, with residue field $\F_2$. The proofs depend only on Lean's three standard
axioms, and neither the SAT solver nor \texttt{drat-trim} is trusted. For the other fields the LRAT proofs,
from $0.7$ to $70$~MB, are too large for \texttt{lrat\_proof} in reasonable time and memory. For them the
Lean files prove the theorem (for $d=47$ the lower bound) from the hypothesis that $\Phi(G)$ is
unsatisfiable: there $\Phi(G)$ is a Lean object, which each file compares with the stored formula by evaluation
(\texttt{\#guard}) when it is built, and \texttt{cake\_lpr} checked an LRAT proof that the stored formula is
unsatisfiable. So for every graph the kernel checks the unit distances, the upper bound (for $d\ne47$) and
the step from the unsatisfiability of $\Phi(G)$ to the theorem, and that unsatisfiability is checked by the
kernel (ten fields) or by \texttt{cake\_lpr}.

\section{The upper bounds}\label{sec:upper}

\begin{lemma}[Moorhouse {\cite[Lemma~8.2]{Moorhouse}}]
If $p\equiv3\pmod4$ is prime and $d$ is $0$ or a nonzero square modulo $p$, then
$\chi(\Q(\sqrt d)^2)\le\chi(\F_p^2)$, where $\F_p^2$ carries the graph $x^2+y^2=1$.
\end{lemma}

The prime $p$ ramifies or splits in $K$; at a prime $\p$ of norm $p$ the form $x^2+y^2$ is anisotropic over
$\F_p$, so every unit vector is $\p$-integral and reduction modulo $\p$ is a graph homomorphism on the
component of the origin; every component is a translate of that one \cite[Lemma~4.2]{Moorhouse}. Since
$\chi(\F_7^2)=4$ \cite[Table~6.1]{Moorhouse} and $\F_{11}^2$ has a proper
$5$-colouring \cite[Lemma~4.5]{Madore}:
\begin{itemize}
\item $d=11$, 23, 71, 95, 155, 179, 191, 239, 263, 359, 431, 443, 491, 599, 611, 851, 911, 935 are nonzero squares
  modulo $7$ and
  $d=35$, 119, 455, 791, 959 are $\equiv0$ modulo $7$, so $\chi\le4$;
\item $d=59,131,251$ are $\equiv3\pmod8$, so $\chi\le4$ by Fischer's Theorem~10 (which also covers $11$, $35$,
  $155$, $179$, $443$, $491$, $611$ and $851$);
\item $47\equiv5^2\pmod{11}$, so $\chi(\Q(\sqrt{47})^2)\le5$.
\end{itemize}
The repository stores a proper $4$-colouring of $\F_7^2$ and a proper $5$-colouring of $\F_{11}^2$, and the
checker checks them. For $d=47$ the next level at $11$ does not help: the unit-distance graph on
$(\Z/121)^2$ is not $4$-colourable.

\begin{corollary}\label{cor:larger}
Let $L\subset\R$ be a number field that contains $\sqrt d$ for one of the $d$ with $\chi(\Q(\sqrt d)^2)=4$ in
Theorem~\ref{thm:main} and has a place with residue field $\F_7$. Then $\chi(L^2)=4$. For example
$L=\Q(\sqrt{11},7^{1/m})$, of degree $2m$; so there are real number fields of every even degree with
$\chi(L^2)=4$.
\end{corollary}

\begin{proof}
The graph over $\Q(\sqrt d)$ lies in $L^2$, so $\chi(L^2)\ge4$. Since $-1$ is not a square in $\F_7$, unit
vectors are integral at the place, and reduction gives $\chi(L^2)\le\chi(\F_7^2)=4$ as above. The prime $7$
splits in $\Q(\sqrt{11})$, and $x^m-7$ is Eisenstein at each prime above it; so $[L:\Q]=2m$, and these
primes are totally ramified in $L$, with residue field $\F_7$.
\end{proof}

By \cite{FC}, the same holds for every such $L$ that contains $\sqrt d$ for some $d\equiv11\pmod{12}$ (added
4 October). Fields of odd degree have $\chi(L^2)=2$ \cite[Theorem~7.1]{Moorhouse}.

\section{How the graphs were found}\label{sec:search}

\subsection*{Directions}
For a denominator $D$ let $U_D$ be the set of unit vectors $[a,b,c,e]$ with integer coordinates, that is, of
solutions of \eqref{eq:unit} with $\Delta=(a,b,c,e)$. It is finite, and closed under negation, the quarter
turn and complex conjugation. We used sets with $68$ to $276$ elements.

\subsection*{Colouring-guided growth}
Start from all sums of at most two vectors of $U_D$; $3$-colour the graph (tabu search, else kissat); add the
candidate points $p+u$ ($p$ in the graph, $u\in U_D$) whose neighbours in the graph already see all three
colours; repeat until kissat answers that the graph is not $3$-colourable. When no candidate was blocked, a
variant coloured the graph again from scratch and, failing that, added candidates whose neighbours see two
colours. The successful runs took $1$ to $96$ rounds and ended with $2\,930$ to $40\,198$ points.

\subsection*{Shrinking}
A single incremental SAT solver (CaDiCaL through PySAT) holds the formula with one selector literal per vertex,
enabling its at-least-one clause. Each test is a solve under assumptions, and each refutation returns a core of
selectors, a subset of the vertices that is still not $3$-colourable. Deleting vertices one at a time, keeping
a deletion when the rest is still refuted, gives a vertex-critical graph; repeating with random deletion orders
(up to $400$, fewer for the largest graphs) and keeping the smallest result gave the graphs of
Table~\ref{tab:graphs}, except four. For $d=191$, $239$, $455$ and $935$ we first kept the vertices whose
clauses lie in the unsatisfiable core that \texttt{drat-trim} extracts from kissat's proof, in rounds, and
started up to $3\,000$ random orders from one of these cores ($122$, $449$, $78$ and $315$ vertices): they
gave $96$, $338$, $71$ and $252$ vertices, where the orders from the grown graphs had given $100$, $355$, $74$
and $257$.

\subsection*{The graph over \texorpdfstring{$\Q(\sqrt{11})$}{Q(sqrt11)}}
With $D=30$ the growth stops at once: the $5\,941$ points within two steps of $0$ have no proper $3$-colouring.
The graph of Figure~\ref{fig:q11} came from shrinking a larger graph: the $11\,612$ points within two steps of
$0$ or of $m=[-15,0,24,-9]$, together with their copy rotated about $0$ by $u=(t+i)/(t-i)$,
$t=(8-3\sqrt{11})/5$. Since $4|m|^2-1=t^2$, this rotation moves $m$ by exactly $1$; we built the union to test a
spindle, as in Moser's, where $|m|^2=3$ and $t=\sqrt{11}$, but it is not needed. Shrinking kept $76$ vertices, all
in the copy; rotated back by $\bar u$, $75$ of them are within two steps of $0$.

\subsection*{Gates}
Some searches stop with $3$-colourable graphs for an arithmetic reason. Let $\p$ be a prime of $K$ above $p$
that splits as $\PP\bar\PP$ in $K(i)$, where the bar is the automorphism $i\mapsto-i$ over $K$, and let
$\F_q$ be its residue field. Writing a unit vector as $z=x+iy$, we have $z\bar z=1$, so
$v_\PP(z)=-v_{\bar\PP}(z)$; if $z$ is integral at $\PP$ and $\bar\PP$ (for instance if $p\nmid D$), its
residues there are inverse to each other. Hence, if every direction is integral at $\PP$ and $\bar\PP$, the
reduction $z\mapsto(z\bmod\PP,\ z\bmod\bar\PP)$ maps every graph built from the directions
homomorphically into
\[
  H_q=\mathrm{Cay}\bigl(\F_q^2,\ \{(x,x^{-1}):x\in\F_q^\times\}\bigr),
\]
and $\chi(H_q)=3$ for $q=3,5,9$, while $\chi(H_q)\ge5$ for $q=13,17,25,29,37,41,49$ (computed with kissat).
This gives three \emph{gates}, each closed whenever $p\nmid D$:
\begin{itemize}
\item at $2$, when $d\equiv7\pmod8$: the prime above $2$ splits in $K(i)$ with $q=2$, and $H_2$ is a perfect
  matching, so every graph built from the directions is bipartite;
\item at $3$, since $d\equiv2\pmod3$: $q=9$. A divisor $3$ of $D$ is not always enough: for $d=83$ and
  $D=1230=2\cdot3\cdot5\cdot41$ no direction has $3$ in its reduced denominator, so every direction is still
  integral at the primes above $3$;
\item at $5$, when $d\equiv0,\pm1\pmod5$: $q=5$.
\end{itemize}
So $D$ must be divisible by $3$, by $2$ when $d\equiv7\pmod8$, and by $5$ when $d\equiv0,\pm1\pmod5$;
every denominator in Table~\ref{tab:graphs} is. A useful set contains vectors that are integral at no gate, such as $(1+7i\sqrt{47})/48$, from Pell's equation $48^2-47\cdot7^2=1$. The gates do not
explain every failure: the first searches over $\Q(\sqrt{35})$, $\Q(\sqrt{131})$ and $\Q(\sqrt{179})$ failed
with sets of $60$ to $212$ directions that pass all three gates, and succeeded with the $196$ to $276$
directions of denominator dividing $390$ and the variant of the growth.

\section{Remarks}

\subsection*{A different question: \texorpdfstring{$\Q(\sqrt{-d})$ inside $\C$}{Q(sqrt(-d)) inside C}}
Cohen \cite{Cohen} wrote that ``the natural conjecture is that $\Q[\alpha]$ is $3$-colorable for all $\alpha$
quadratic over $\Q$'', for the set $\Q[\alpha]\subset\C$. For $\alpha=\sqrt{-d}$ this set is
$\{x+y\sqrt{-d}\}$, of dimension $2$ over $\Q$, while $\Q(\sqrt d)^2$ has dimension $4$; so
Theorem~\ref{thm:main} does not contradict the conjecture. The conjecture holds:

\begin{proposition}
For every $\alpha$ quadratic over $\Q$, the unit-distance graph on $\Q[\alpha]\subset\C$ is $3$-colourable.
\end{proposition}

\begin{proof}
For real $\alpha$ the set lies on a line. For $\alpha=\sqrt{-1}$ it is $\Q^2$ \cite{Woodall}, and for
$\alpha=\sqrt{-2}$ it lies in $\Q(\sqrt2)^2$, which is $2$-colourable. Otherwise $d$ has an odd prime factor
$p$, and the prime $\p$ of $E=\Q(\sqrt{-d})$ above $p$ is ramified. An element $z\in E$ with $|z|=1$
satisfies $z\bar z=1$ and $v_\p(\bar z)=v_\p(z)$, so $z$ is a $\p$-adic unit; and $z\equiv\bar z\pmod\p$, so $z^2\equiv1$ and
$z\equiv\pm1\pmod\p$. Choose a representative $r_C$ of each coset $C$ of the ring $\mathcal O_\p$ of
$\p$-integral elements, and colour $w\in C$ by $f((w-r_C)\bmod\p)$, where $f$ is a proper $3$-colouring of the
cycle $\mathrm{Cay}(\F_p,\{\pm1\})$. Points at distance $1$ differ by such an element, lie in the same coset,
and have residues differing by $\pm1$.
\end{proof}

\subsection*{Open questions}
\begin{itemize}
\item Is $\chi(\Q(\sqrt{47})^2)$ equal to $4$ or $5$? A value of $5$ would give a triangle-free $5$-chromatic
  unit-distance graph in the plane. The bound $5$ is the reduction at a prime above $11$, where $\Q(\sqrt{47})$
  embeds in $\Q_{11}$; every level $\mathrm{Cay}((\Z/11^k)^2,\{x^2+y^2=1\})$ of it bounds $\chi$, but the levels
  $k=1,2,3$ have no proper $4$-colouring: a $69$-point set of level $1$ with none lifts to a $244$-point set of
  level $2$ with none, and the $29\,524$ points of level $3$ above it have none either (kissat and drat-trim). The same $11$-adic plane bounds $\chi(\Q(\sqrt3,\sqrt5)^2)$, since $\Q(\sqrt3,\sqrt5)$
  embeds in $\Q_{11}$ as well. No other place does better: where $-1$ is a square in $K_v$, no colouring of
  $K_v^2$ with finitely many colours is locally constant, and above a prime $p\equiv3\pmod4$ with $p\ge23$,
  Hoffman's bound gives every level an independence ratio below $1/4$ (at level $1$ the exact least eigenvalue
  gives the ratios $0.2491$ and $0.249997$ for $p=23$ and $31$, and for $p\ge43$ the bound $2\sqrt p$ on the
  eigenvalues \cite{MMST} suffices; the eigenvalues new at higher levels are at most $2/(p+1)$ of the degree), and
  so does the ramified place above $47$. So a $4$-colouring by reduction would come from the $11$-adic plane at level $4$ or
  more, or from the $19$-adic plane at level $2$ or more. (Added 4 October: by \cite{TCol},
  $4\le\chi_c(\Q(\sqrt{47})^2)\le19/4$.)
\item Does every real quadratic field with $d\equiv11\pmod{12}$ need four colours? Our searches over
  $\Q(\sqrt{83})$, $\Q(\sqrt{107})$ and $\Q(\sqrt{203})$ found only $3$-colourable graphs. (Answered since,
  4 October: yes \cite{FC}; for these three fields $\chi=4$ \cite{Winding}. No explicit graph is known for them.)
\item A uniform proof: a family of graphs, or an argument, covering all $d\equiv11\pmod{12}$ at once. (The argument
  now exists \cite{FC}, and is not constructive; a uniform family of explicit graphs does not.)
\item How small can a $4$-chromatic unit-distance graph over a real quadratic field be? The smallest of ours
  has $71$ vertices, for $d=455$; the one for $d=11$ has $76$.
\end{itemize}

\begin{thebibliography}{99}
\bibitem{cakelpr} Y. K. Tan, M. J. H. Heule and M. O. Myreen, \emph{cake\_lpr: verified propagation redundancy
checking in CakeML}, TACAS 2021, Lecture Notes in Comput. Sci. \textbf{12652}, 223--241.
\bibitem{Cohen} D. Cohen, \emph{The $\C$ unit distance graph}, University of Chicago REU paper, 2007,
\url{https://www.math.uchicago.edu/~may/VIGRE/VIGRE2007/REUPapers/FINALFULL/Cohen.pdf}.
\bibitem{deGrey} A. D. N. J. de Grey, \emph{The chromatic number of the plane is at least 5}, Geombinatorics
\textbf{28} (2018), 18--31; arXiv:1804.02385.
\bibitem{Fischer1990} K. G. Fischer, \emph{Additive $K$-colorable extensions of the rational plane}, Discrete
Math. \textbf{82} (1990), 181--195.
\bibitem{Repo} S. Gal\'an, \emph{The Hadwiger--Nelson problem over number fields}, code, data and formal
proofs, 2026. Archived at Zenodo, \href{https://doi.org/10.5281/zenodo.22976635}{doi:10.5281/zenodo.22976635}
(all versions); source at \url{https://github.com/decalion89/chromatic-number-of-the-plane}; a fuller account is
\texttt{notes/quadratic\_planes.md} there. These files and the Lean files cited here are on the branch
\texttt{claude/quadratic-planes} (pull request~35), not yet in a release.
\bibitem{FC} S. Gal\'an, \emph{Which real quadratic planes need four colours}, draft, 2026,
\texttt{papers/four-colours/} in \cite{Repo}.
\bibitem{TCol} S. Gal\'an, \emph{Three colours for the plane over a number field}, draft, 2026,
\texttt{papers/three-colours/} in \cite{Repo}.
\bibitem{Winding} S. Gal\'an, \emph{Circular colourings of abelian Cayley graphs below four come from
characters}, draft, 2026, \texttt{papers/winding/} in \cite{Repo}.
\bibitem{Johnson1987} P. D. Johnson Jr., \emph{Two-colorings of real quadratic extensions of $\Q^2$ that forbid
many distances}, Congr. Numer. \textbf{60} (1987), 51--58.
\bibitem{Madore} D. A. Madore, \emph{The Hadwiger--Nelson problem over certain fields}, arXiv:1509.07023 (2015).
\bibitem{mathlib} The mathlib Community, \emph{The Lean mathematical library}, CPP 2020, 367--381;
arXiv:1910.09336.
\bibitem{Lean} L. de Moura and S. Ullrich, \emph{The Lean 4 theorem prover and programming language}, CADE 28,
Lecture Notes in Comput. Sci. \textbf{12699} (2021), 625--635.
\bibitem{MMST} A. Medrano, P. Myers, H. M. Stark and A. Terras, \emph{Finite analogues of Euclidean space},
  J. Comput. Appl. Math. \textbf{68} (1996), 221--238.
\bibitem{MM} MildlyMeticulous (GitHub account), \emph{A $2$-adic obstruction to $5$-chromatic unit-distance
graphs}, code, logs and draft, July 2026,
\url{https://github.com/MildlyMeticulous/hn-2adic-obstruction}.
\bibitem{Moorhouse} G. E. Moorhouse, \emph{On the chromatic numbers of planes}, preliminary draft, revised
3 March 2010, \url{https://www.ericmoorhouse.org/pub/chromatic.pdf}.
\bibitem{MoorhouseTalk} G. E. Moorhouse, \emph{Colouring the plane}, talk at Designs, Codes \& Geometries
2010, slides at \url{https://www.ericmoorhouse.org/slides/ebertfest.pdf}.
\bibitem{Payne} M. S. Payne, \emph{Unit distance graphs with ambiguous chromatic number}, Electron. J. Combin.
\textbf{16} (2009), Note~31; arXiv:0707.1177.
\bibitem{Woodall} D. R. Woodall, \emph{Distances realized by sets covering the plane}, J. Combin. Theory Ser. A
\textbf{14} (1973), 187--200.
\end{thebibliography}

\end{document}
